\section{Related Work}

Compile-time metaprogramming in ISO C++ has progressed from template tricks and macro preprocessors to standardized \texttt{constexpr} facilities \cite{p1240}, culminating in P2996 \cite{p2996} and P1967 \cite{p1967}. CTRE \cite{ctre} established the precedent that an entire domain-specific grammar, in that case regular expressions, can be tokenized and compiled to typed C++ artifacts inside \texttt{constexpr} evaluation rather than through an external tool; \texttt{ctpb} applies the same principle to the substantially larger proto3 grammar and adds type synthesis rather than pattern matching. Boost.PFR \cite{boost_pfr} predates P2996 and offers structural reflection over existing aggregates, sufficient for serializing a struct that already exists but unable to synthesize members from an external schema.

Among serialization libraries, reflect-cpp \cite{reflect_cpp} provides C++20 reflection-based serialization but derives schemas from existing C++ structs rather than compiling external IDL, and so lacks proto3 wire tags, packed encoding, and reserved-range validation. FlatBuffers \cite{flatbuffers} and Cap'n Proto \cite{capnproto} target zero-copy in-memory layout but ship standalone compiler binaries analogous to \texttt{protoc}. The lightweight protobuf runtimes reduce overhead without removing the generation step: \texttt{libprotobuf-lite} omits runtime reflection descriptors, nanopb \cite{nanopb} targets embedded systems through its own \texttt{protoc} plugin, protozero \cite{protozero} relies on hand-written accessors, and \texttt{upb} \cite{upb}, Google's own small-footprint C core, still consumes \texttt{protoc}-generated descriptors. \texttt{ctpb} differs from all of these in eliminating the generation step outright while retaining proto3 as the schema authority.
